% ============================================================
%  Chapitre 7 — Discussion et conclusion
% ============================================================
\chapter{Discussion et conclusion}
\label{chap:discussion}

\section{Synthèse comparative des performances}

Le tableau~\ref{tab:synth} regroupe les meilleures configurations obtenues pour
chaque approche de parallélisation.

\begin{table}[h]
\centering
\caption{Synthèse des performances — toutes versions}
\label{tab:synth}
\begin{tabular}{llrrr}
\hline
\textbf{Version}              & \textbf{Config.}    & \textbf{$T$ (ms/it)} & \textbf{$S$} & \textbf{$E$ (\%)} \\
\hline
AoS séquentiel (ref.)         & 1 thread            & 1.75  & 1.0 & 100 \\
SoA + OpenMP                  & 1 thread            & 2.09  & 0.8 & --- \\
SoA + OpenMP                  & 4 threads           & 0.66  & 2.7 & 66 \\
SoA + OpenMP                  & 14 threads          & 0.29  & 6.0 & 43 \\
MPI Approche 1                & 1 processus         & 1.57  & 1.1 & --- \\
MPI Approche 1                & 2 processus         & 1.51  & 1.2 & 58 \\
\hline
\end{tabular}
\end{table}

\section{Discussion}

\subsection{OpenMP : stratégie gagnante pour ce problème}

OpenMP est clairement l'approche la plus efficace pour ce cas d'usage, avec une
accélération de $7.1\times$ à 14 threads. La nature \emph{embarrassingly parallel}
de la boucle sur les fourmis et la bonne localité mémoire de la SoA favorisent
cette approche.

La saturation observée à partir de 8 threads s'explique par deux facteurs concurrents :
la fraction séquentielle résiduelle ($\approx 7\%$, Amdahl) et la limite de bande
passante mémoire pour l'évaporation. La loi de Gustafson prédit qu'un passage à
$m = 50\,000$ fourmis permettrait de maintenir une meilleure efficacité à grand
nombre de threads (fraction parallèle augmentée), mais au prix d'une consommation
mémoire plus importante.

\subsection{MPI Approche 1 : bottleneck de communication}

L'Approche 1 MPI est inadaptée à ce problème dans sa configuration actuelle.
Le ratio \emph{calcul/communication} est défavorable : chaque processus ne fait que
$m/P \approx 2\,500$ appels à \texttt{advance\_one()} par itération, mais synchronise
systématiquement $4.2$~Mo de phéromones. Ce déséquilibre est intrinsèque à la stratégie
de carte globale et ne peut pas être résolu sans réduire $N$ ou augmenter $m$.

\subsection{MPI Approche 2 : voie prometteuse mais complexe}

L'Approche 2 (décomposition de domaine) réduit la communication de $260\times$ et
permettrait théoriquement une bonne scalabilité. Cependant, sa complexité d'implémentation
est nettement supérieure, notamment la gestion de la migration des fourmis et
l'équilibrage de charge dynamique. Elle est pertinente pour des grilles très grandes
($N \geq 4\,096$) ou pour des environnements à mémoire distribuée où l'Approche 1
est impossible.

\subsection{Limites de l'étude}

Plusieurs facteurs limitent la portée des conclusions :

\begin{itemize}
    \item Les mesures ont été effectuées sur un seul nœud de calcul (machine locale),
          sans réseau inter-nœuds pour MPI. Les latences réelles d'un cluster HPC
          seraient différentes.
    \item La race condition dans \texttt{mark\_pheronome()} rend les résultats
          fonctionnels non-déterministes en OpenMP. Bien que les phéromones
          convergent statistiquement, l'indicateur ``nourriture collectée'' varie
          légèrement entre runs (jusqu'à $\pm 5\%$).
    \item La comparaison AoS vs SoA est bruitée par le surcoût du runtime OpenMP
          à 1 thread. Une comparaison équitable nécessiterait de compiler les deux
          versions sans \texttt{-fopenmp}.
\end{itemize}

\section{Conclusion}

Ce projet a permis d'explorer progressivement les trois niveaux de parallélisation
d'un algorithme ACO sur terrain fractal : instrumentation et profiling, vectorisation
SoA, parallélisation OpenMP en mémoire partagée, et parallélisation MPI distribuée.

Les principaux enseignements sont :

\begin{enumerate}
    \item \textbf{Profiler avant d'optimiser} : la phase fourmis ($90\%$ du temps) est
          le seul goulot d'étranglement réel. Optimiser l'évaporation ($9\%$) au-delà
          d'un certain point n'apporte pas de gain global significatif.
    \item \textbf{OpenMP est adapté} aux algorithmes d'agents indépendants en mémoire
          partagée : $7\times$ de speedup à 14 threads avec un effort d'implémentation
          limité.
    \item \textbf{MPI Approche 1 est limitée par la communication} : le volume de
          données synchronisé ($4.2$~Mo/it) annule le gain de calcul dès 2 processus.
          L'Approche 2 serait plus efficace mais nettement plus complexe à réaliser.
    \item \textbf{L'architecture hybride moderne} (P-cores + E-cores) crée des
          effets de saturation non linéaires qu'il faut identifier empiriquement.
\end{enumerate}

En perspective, une combinaison \textbf{MPI + OpenMP} (modèle hybride) semblerait
prometteuse : MPI Approche 2 pour la décomposition de domaine entre nœuds, OpenMP
pour la parallélisation intra-nœud de chaque bande. Cette architecture exploiterait
pleinement les ressources d'un cluster HPC moderne.

% TODO : ajouter comparaison avec résultats théoriques ACO (convergence, qualité solution)
